% STAGE08-CHAPTER: V1-C05
% CHAPTER-SCHEMA-COVERAGE: CH-01,CH-02,CH-03,CH-04,CH-05,CH-06,CH-07,CH-08,CH-09,CH-10,CH-11,CH-12,CH-13,CH-14,CH-15,CH-16,CH-17,CH-18,CH-19,CH-20,CH-21,CH-22,CH-23,CH-24,CH-25,CH-26,CH-27,CH-28,CH-29,CH-30,CH-31,CH-32,CH-33,CH-34,CH-35,CH-36,CH-37
% CHAPTER-SCHEMA-EVIDENCE: CH-01=\label{struct:V1-C05-CH01}; CH-02=\label{struct:V1-C05-CH02}; CH-03=\label{struct:V1-C05-CH03}; CH-04=\label{struct:V1-C05-CH04}; CH-05=\label{struct:V1-C05-CH05}; CH-06=\label{struct:V1-C05-CH06}; CH-07=\label{struct:V1-C05-CH07}; CH-08=\label{struct:V1-C05-CH08}; CH-09=\label{struct:V1-C05-CH09}; CH-10=\label{struct:V1-C05-CH10}; CH-11=\label{struct:V1-C05-CH11}; CH-12=\label{struct:V1-C05-CH12}; CH-13=\label{struct:V1-C05-CH13}; CH-14=\label{struct:V1-C05-CH14}; CH-15=\label{struct:V1-C05-CH15}; CH-16=\label{struct:V1-C05-CH16}; CH-17=\label{struct:V1-C05-CH17}; CH-18=\label{struct:V1-C05-CH18}; CH-19=\label{struct:V1-C05-CH19}; CH-20=\label{struct:V1-C05-CH20}; CH-21=\label{struct:V1-C05-CH21}; CH-22=\label{struct:V1-C05-CH22}; CH-23=\label{struct:V1-C05-CH23}; CH-24=\label{struct:V1-C05-CH24}; CH-25=\label{struct:V1-C05-CH25}; CH-26=\label{struct:V1-C05-CH26}; CH-27=\label{struct:V1-C05-CH27}; CH-28=\label{struct:V1-C05-CH28}; CH-29=\label{struct:V1-C05-CH29}; CH-30=\label{struct:V1-C05-CH30}; CH-31=\label{struct:V1-C05-CH31}; CH-32=\label{struct:V1-C05-CH32}; CH-33=\label{struct:V1-C05-CH33}; CH-34=\label{struct:V1-C05-CH34}; CH-35=\label{struct:V1-C05-CH35}; CH-36=\label{struct:V1-C05-CH36}; CH-37=\label{struct:V1-C05-CH37}

% FIGURE-KNOWLEDGE-MAP: fig:V1-C05-gaussian -> PRE-PR-019
\chapter{常用分布与统计推断}\label{chap:V1-C05}
\SLChapterOpening{从常用分布直达似然、后验与指数族}{样本怎样通过概率模型约束未知参数}{能够处理支持边界、共轭更新和配分函数导数}{\cref{chap:V1-C04}的概率、条件分布与矩}{支持或参数化不清楚时返回分布定义与似然节}
\index{统计推断}


\phantomsection\label{struct:V1-C05-CH01}
\chaptergoals{
\begin{itemize}[leftmargin=2em]
  \item 掌握离散计数分布、高斯分布、Gamma分布与狄利克雷分布的参数和支持集；
  \item 从抽样模型完整推导极大似然、最大后验与共轭更新，不混淆概率和似然；
  \item 能够处理边缘似然、隐变量与指数族表示，并判断估计方法的适用条件。
\end{itemize}}

\phantomsection\label{struct:V1-C05-CH02}
\begin{prerequisitebox}
本章依赖第四章的条件概率、期望、方差与独立同分布抽样，还使用求和、积分、对数、偏导数和正定矩阵。多元高斯的二次型可回看第二章。
\end{prerequisitebox}

\phantomsection\label{struct:V1-C05-CH03}
% KNOWLEDGE-PLAN: KN-V1-C05-PREREQUISITE-004 L1
\paragraph{读前自检：先修自检。}常用分布与统计推断先修自检固定核验“能否说明概率质量函数对随机结果归一化，而似然函数把观测固定后比较参数”；请说明概率质量函数对随机结果归一化，而似然函数把观测固定后比较参数并写出中间结果，再按“中间结果直接回答首项“能否说明概率质量函数对随机结果归一化，而似然函数把观测固定后比较参数”并给出通过或失败”明确判为通过或失败。\par

\begin{precheckbox}
\begin{enumerate}[leftmargin=2.2em]
  \item 能否说明概率质量函数对随机结果归一化，而似然函数把观测固定后比较参数？不能则回看\GlobalChapterRef{4}{1}{4}的\cref{sec:V1-C04-S04}，并预习本章\cref{sec:V1-C05-S05}。
  \item 能否计算 $\sum_i\log p(x_i\mid\theta)$ 对参数的导数？不能则回看\GlobalChapterRef{3}{1}{3}的\cref{sec:V1-C03-S02,sec:V1-C03-S04,sec:V1-C03-S05,sec:V1-C03-S06}。
  \item 能否判断协方差矩阵为何必须对称半正定？不能则回看\GlobalChapterRef{2}{1}{2}的\cref{sec:V1-C02-S06}或\GlobalChapterRef{4}{1}{4}的\cref{sec:V1-C04-S03}。
\end{enumerate}
\end{precheckbox}

\phantomsection\label{struct:V1-C05-CH04}
\begin{dependencybox}
概率公理 $\to$ 常用分布 $\to$ 独立抽样 $\to$ 似然与极大似然；
先验 $+$ 似然 $\to$ 后验与MAP $\to$ 共轭更新；
联合模型 $\to$ 边缘似然 $\to$ 隐变量推断与指数族统一表示。
\end{dependencybox}

\phantomsection\label{struct:V1-C05-CH05}
% STAGE19-SYMBOL-INDEX-BEGIN: V1-C05
\index[symbols]{minus-theta-minus-minus-theta-minus--sb707ad2907cc@$\theta,\Theta$：未知但固定的模型参数及其可行集合}
\index[symbols]{x-1-u003a-n--s7e30735647f5@$x_{1:n}$：$x_1,\ldots,x_n$ 的简写}
\index[symbols]{l-minus-theta-minus-minus-ell-minus-minus-theta-minus--s3f5b2fa475df@$L(\theta),\ell(\theta)$：似然与对数似然}
\index[symbols]{minus-mu-minus-minus-sigma-minus-2-minus-sigma-minus--s7f033a7082f1@$\mu,\sigma^2,\Sigma$：均值、方差和协方差}
\index[symbols]{minus-alpha-minus-minus-boldsymbol-minus-minus-alpha-minus--s10e0ba694086@$\alpha,\boldsymbol\alpha,\alpha_0$：标量形状参数、Dirichlet参数向量及总浓度}
% STAGE19-SYMBOL-INDEX-END: V1-C05
% KNOWLEDGE-PLAN: KN-V1-C05-PREREQUISITE-006 L1
\paragraph{读前自检：本章符号表。}常用分布与统计推断符号表规定$\theta,\Theta$、$x_{1:n}$的类型或范围；请把$\theta,\Theta$代入紧随其后的定义关系，逐项核对输入形状、输出形状与取值域。\par

\begin{symboltablebox}
\begin{tabularx}{\linewidth}{@{}l l X@{}}
\toprule 符号 & 取值空间 & 含义\\ \midrule
$\theta,\Theta$ & 参数、参数空间 & 未知但固定的模型参数及其可行集合\\
$x_{1:n}$ & 观测样本 & $x_1,\ldots,x_n$ 的简写\\
$L(\theta),\ell(\theta)$ & 非负函数、实函数 & 似然与对数似然\\
$\mu,\sigma^2,\Sigma$ & 实数、正数、对称半正定矩阵 & 均值、方差和协方差；高斯密度式另要求$\Sigma\succ0$\\
$\mathbf N,\mathbf n$ & $\mathbb Z_{\geq0}^K$中的随机向量、实现值 & $K$类计数向量及其一次实现\\
$n$ & 正整数 & 多项分布的总试验数，满足$\sum_{k=1}^K n_k=n$\\
$\alpha$；$\boldsymbol\alpha$；$\alpha_0$ & 正数；$\R_{>0}^K$；正数 & 标量形状参数；Dirichlet参数向量；总浓度$\alpha_0=\sum_k\alpha_k$\\
\bottomrule
\end{tabularx}
\end{symboltablebox}

% KNOWLEDGE-PLAN: KN-V1-C05-CONCEPT-001 L1
\paragraph{读前自检：伯努利分布。}把伯努利分布放在“$X\in\{0,1\}$ 且 $P(X=x)=p^x(1-p)^{1-x}$”下考察：把$P(X=x)=p^x(1-p)^{1-x}$在其支持上求和或积分一次，并检查其期望为 $p$，方差为 $p(1-p)$。\par

\SLDirectSection{伯努利、二项与多项分布}{sec:V1-C05-S01}
\phantomsection\label{struct:V1-C05-CH06}
\knowledgeanchor{PRE-PR-016}{伯努利分布}
\begin{definition}[伯努利分布]\label{struct:V1-C05-CH07}
$X\in\{0,1\}$ 且 $P(X=x)=p^x(1-p)^{1-x}$，其中 $0\le p\le1$。其期望为 $p$，方差为 $p(1-p)$。
\end{definition}
\knowledgeanchor{PRE-PR-017}{二项分布}
若 $X_1,\ldots,X_n$ 独立且同服从 $\operatorname{Bernoulli}(p)$，成功次数 $S=\sum_iX_i$ 的质量函数为
\[
P(S=k)=\binom nk p^k(1-p)^{n-k},\quad k=0,\ldots,n.
\]
\knowledgeanchor{PRE-PR-018}{多项分布}
把每次试验推广到$K$个类别，记随机计数向量与其实现值分别为
$\mathbf N=(N_1,\ldots,N_K)$和$\mathbf n=(n_1,\ldots,n_K)$；标量$n$专指总试验数，
故$\sum_{k=1}^K N_k=n$且可行实现满足$\sum_{k=1}^K n_k=n$。此时
\[
P(\mathbf N=\mathbf n)=\frac{n!}{\prod_{k=1}^K n_k!}
\prod_{k=1}^{K}p_k^{n_k},\qquad
p_k\geq0,\quad\sum_{k=1}^Kp_k=1.
\]

\phantomsection\label{struct:V1-C05-CH12}
\begin{theorem}[二项分布的均值与方差]
若 $S\sim\operatorname{Binomial}(n,p)$，则 $\mathbb ES=np$，$\operatorname{Var}(S)=np(1-p)$。
\end{theorem}
\phantomsection\label{struct:V1-C05-CH15}
% KNOWLEDGE-PLAN: KN-V1-C05-PROOF-001 L1
\paragraph{读前自检：证明：二项分布的均值与方差。}二项分布的均值与方差要求先确认“$\sum_i p(1-p)=np(1-p)$”；请随后执行$\sum_i p(1-p)=np(1-p)$中最内层的一次求和、积分或期望，用独立性使协方差项为零，故和的方差为 $\sum_i p(1-p)=np(1-p)$验收。\par
% KNOWLEDGE-PLAN: KN-V1-C05-CONCEPT-004 L1
\paragraph{读前自检：一元高斯分布。}已知$\sigma^2>0$。请说明高斯密度为何非负，并核对其在$\mathbb R$上的积分为$1$。\par

\begin{proof}
把 $S$ 写成独立伯努利变量之和。期望线性性给出 $\sum_i p=np$；独立性使协方差项为零，故和的方差为 $\sum_i p(1-p)=np(1-p)$。
\end{proof}
\SLReviewSection{高斯分布与多元高斯分布}{sec:V1-C05-S02}
\knowledgeanchor{PRE-PR-019}{一元高斯分布}
一元高斯密度为
\[
f(x\mid\mu,\sigma^2)=\frac{1}{\sqrt{2\pi\sigma^2}}
\exp\!\left[-\frac{(x-\mu)^2}{2\sigma^2}\right],\qquad \sigma^2>0.
\]
\knowledgeanchor{PRE-PR-033}{多元高斯分布}
对$x\in\mathbb R^d$，当$\Sigma\succ0$时，多元高斯是非退化分布，并具有相对于
$d$维Lebesgue测度的密度
\[
f(x)=\frac{\exp[-\tfrac12(x-\mu)^{\mathsf T}\Sigma^{-1}(x-\mu)]}
{(2\pi)^{d/2}|\Sigma|^{1/2}}.
\]
二次型给出马氏距离，行列式控制等密度椭球的体积尺度。

\begin{keypointbox}[title=扩展讨论：非退化密度与退化高斯测度]
若只要求$\Sigma\succeq0$，仍可通过$X=\mu+LZ$、$LL^{\mathsf T}=\Sigma$定义高斯测度；
当$\operatorname{rank}(\Sigma)<d$时，它集中在低维仿射子空间
$\mu+\operatorname{col}(L)$上，相对于$d$维Lebesgue测度没有密度，因而不能把
$|\Sigma|=0$或不存在的$\Sigma^{-1}$代入上式。

一个秩亏二维反例是：令$Z\sim\mathcal N(0,1)$并取$X=(Z,Z)^{\mathsf T}$，则
\[
\Sigma=\begin{bmatrix}1&1\\1&1\end{bmatrix}\succeq0,
\qquad \operatorname{rank}(\Sigma)=1.
\]
这是一个退化二维高斯向量，其全部质量集中在直线$x_1=x_2$上；它是高斯测度，
却不是可由上述二维密度公式表示的非退化高斯密度。
\end{keypointbox}

\begin{lemma}[标准化]
若 $X\sim\mathcal N(\mu,\sigma^2)$，则 $Z=(X-\mu)/\sigma\sim\mathcal N(0,1)$。
\end{lemma}
% KNOWLEDGE-PLAN: KN-V1-C05-PROOF-002 L1
\paragraph{读前自检：证明：标准化。}常用分布与统计推断中的标准化依赖“把原密度与雅可比相乘后，$\sigma$ 与归一化常数中的 $1/\sigma$ 抵消”；请据此把$x=\mu+\sigma z$用到的关键定义展开并完成下一步变形，并直接核对得到标准高斯密度。\par

\begin{proof}
使用变量代换 $x=\mu+\sigma z$，雅可比绝对值为 $\sigma$。把原密度与雅可比相乘后，$\sigma$ 与归一化常数中的 $1/\sigma$ 抵消，指数项化为 $-z^2/2$，因此得到标准高斯密度。
\end{proof}

\phantomsection\label{struct:V1-C05-CH19}
% v2.6.0 figure UID: FIG-P077-01
% Image-reference reverse redraw: exact Gaussian formulas retained; comparison hierarchy and label clearance refined.
\tikzset{
  slfig-FIG-P077-01/.style={every path/.append style={line cap=round,line join=round}},
  slfig-FIG-P077-01-label/.style={
    slfig direct label,font=\fontsize{9.2pt}{11.2pt}\selectfont,
    fill=white,fill opacity=.95,text opacity=1,inner xsep=1.1pt,inner ysep=.8pt
  }
}
\pgfplotsset{slfig-FIG-P077-01-axis/.style={
  axis line style={draw=SLInk!90,line width=.60pt},
  tick style={draw=SLInk!80,line width=.55pt},
  tick label style={font=\fontsize{8.8pt}{10.5pt}\selectfont},
  label style={font=\fontsize{9.4pt}{11.2pt}\selectfont},
  clip=false
}}
\pgfkeysifdefined{/tikz/alt}{}{\tikzset{alt/.code={}}}
\begin{figure}[H]
\centering
\begin{tikzpicture}[slfig-FIG-P077-01,alt={对象—关系—结论：方差增大使高斯密度变宽、峰值降低，但曲线下面积保持为1}]
\begin{axis}[slfig-FIG-P077-01-axis,
  slfig axis,width=.86\linewidth,height=6.2cm,
  axis lines=left,xlabel={$x$},ylabel={密度},
  xmin=-4.5,xmax=4.5,ymin=-.065,ymax=.43,
  domain=-4.5:4.5,samples=321,smooth,
  xtick={-4,0,4},ytick={0,.2,.4},clip=false
]
  \path[name path=zero] (axis cs:-4.5,0)--(axis cs:4.5,0);
  \addplot[name path=wide,draw=none,no marks]
    {1/sqrt(8*pi)*exp(-x^2/8)};
  \addplot[fill=SLTeal,fill opacity=.08,draw=none]
    fill between[of=wide and zero];
  \addplot[name path=narrow,draw=none,no marks]
    {1/sqrt(2*pi)*exp(-x^2/2)};
  \addplot[fill=SLBlue,fill opacity=.09,draw=none]
    fill between[of=narrow and zero];
  \addplot[draw=SLBlue,line width=1.05pt,no marks]
    {1/sqrt(2*pi)*exp(-x^2/2)};
  \addplot[draw=SLTeal,line width=.90pt,
    dash pattern=on 3.4pt off 2.0pt,no marks]
    {1/sqrt(8*pi)*exp(-x^2/8)};
  \draw[draw=SLGray!85,line width=.55pt,dash pattern=on 2.8pt off 1.8pt]
    (axis cs:0,0)--(axis cs:0,.405);
  \node[slfig-FIG-P077-01-label,anchor=south west] at (axis cs:.95,.305)
    {$\mathcal N(0,1)$；峰值 $1/\sqrt{2\pi}$};
  \node[slfig-FIG-P077-01-label,anchor=south west] at (axis cs:2.05,.135)
    {$\mathcal N(0,2^2)$；峰值 $1/(2\sqrt{2\pi})$};
  \draw[draw=SLGray!88,line width=.52pt,decorate,
    decoration={brace,mirror,amplitude=3pt}]
    (axis cs:-3.5,-.018)--(axis cs:3.5,-.018)
    node[pos=.50,below=3pt,font=\fontsize{9.2pt}{11.0pt}\selectfont,
      fill=white,inner sep=.8pt]
    {两条曲线：面积 $=1$};
\end{axis}
\end{tikzpicture}
\caption{方差增大使高斯密度变宽、峰值降低，但曲线下面积保持为1}\label{fig:V1-C05-gaussian}
\end{figure}

比较\cref{fig:V1-C05-gaussian}中的曲线时要同时检查位置、尺度和面积：改变$\mu$只平移中心，改变$\sigma$会展宽并降低峰值；若只横向拉伸而不调整$1/(\sqrt{2\pi}\sigma)$，所得曲线便不再是概率密度。


\SLReviewSection{Gamma与狄利克雷分布}{sec:V1-C05-S03}
\knowledgeanchor{PRE-PR-031}{Gamma函数与Gamma分布}
Gamma函数定义为 $\Gamma(a)=\int_0^\infty t^{a-1}e^{-t}\,\mathrm dt$（$a>0$），分部积分给出 $\Gamma(a+1)=a\Gamma(a)$。采用形状--速率参数化时，
\[
f(x\mid a,b)=\frac{b^a}{\Gamma(a)}x^{a-1}e^{-bx},\quad x>0,
\]
其均值为 $a/b$，方差为 $a/b^2$。狄利克雷分布定义在概率单纯形上，后续作为多项分布的共轭先验。

\phantomsection\label{struct:V1-C05-CH13}
% KNOWLEDGE-PLAN: KN-V1-C05-LEMMA-002 L1
\paragraph{读前自检：Gamma递推。}Gamma递推以“特别地，正整数 $n$ 满足 $\Gamma(n)=(n-1)!$”为约束；请独立化简$\Gamma(a+1)=a\Gamma(a)$的两端，并确认两端运算均有定义、类型一致且化为同一结果。\par

\begin{lemma}[Gamma递推]
对 $a>0$，$\Gamma(a+1)=a\Gamma(a)$；特别地，正整数 $n$ 满足 $\Gamma(n)=(n-1)!$。
\end{lemma}
\begin{proof}
对 $\int_0^\infty t^ae^{-t}\,\mathrm dt$ 分部积分，取 $u=t^a$、$\mathrm dv=e^{-t}\mathrm dt$。边界项在零与无穷处均为零，剩余积分为 $a\int_0^\infty t^{a-1}e^{-t}\,\mathrm dt$。再从 $\Gamma(1)=1$ 递推到整数即可。
\end{proof}


\SLTeachSection{总体、样本、参数与统计量}{sec:V1-C05-S04}
\knowledgeanchor{PRE-PR-020}{总体、样本、参数与统计量}
总体分布 $P_\theta$ 是数据生成机制，$\theta$ 是描述该机制的未知参数；样本 $X_{1:n}$ 是从总体抽取的随机变量组，观察后得到固定数值 $x_{1:n}$。统计量 $T(X_{1:n})$ 只能依赖样本，不能依赖未知参数。估计量是用于估计参数的统计量，其抽样分布描述重复抽样时的波动。


% KNOWLEDGE-PLAN: KN-V1-C05-FORMULA-001 L1
\paragraph{读前自检：似然函数。}似然函数在观测$x_{1:n}$固定后把$\theta$视为变量；请在条件独立假设下写出$L(\theta;x_{1:n})=\prod_{i=1}^n p(x_i\mid\theta)$，并核对它无需对$\theta$积分为1但必须尊重参数域。\par
% KNOWLEDGE-PLAN: KN-V1-C05-FORMULA-002 L1
\paragraph{读前自检：对数似然。}对常用分布与统计推断中的对数似然，条件“严格单调的对数把乘积化为和：$\ell(\theta)=\sum_i\log p(x_i\mid\theta)$”不可省；请把观测对应因子代入$\ell(\theta)=\sum_i\log p(x_i\mid\theta)$并合并一次，再用所得目标保留参数定义域且无需对参数归一化作检查。\par

\SLTeachSection{似然、对数似然与极大似然}{sec:V1-C05-S05}
\knowledgeanchor{PRE-PR-021}{似然函数}
给定观测 $x_{1:n}$ 后，$L(\theta;x)=p(x_{1:n}\mid\theta)$ 被视为参数 $\theta$ 的函数。独立抽样时它是单点概率或密度的乘积。
\knowledgeanchor{PRE-PR-022}{对数似然}
严格单调的对数把乘积化为和：$\ell(\theta)=\sum_i\log p(x_i\mid\theta)$，且不改变最大点。
\knowledgeanchor{PRE-PR-023}{极大似然估计}
\[
\hat\theta_{\mathrm{MLE}}\in\arg\max_{\theta\in\Theta}\ell(\theta).
\]

\phantomsection\label{struct:V1-C05-CH08}
\textbf{模型。}伯努利抽样模型假定给定$p$后样本条件独立，数据只通过成功次数进入似然。
\phantomsection\label{struct:V1-C05-CH09}
\textbf{学习策略。}选择使已观测数据密度最大的可行参数，并用边界、二阶条件或凸性核实该驻点是否为全局最大点。

% DERIVATION-CHECK: step-1; step-2; step-3
\phantomsection\label{struct:V1-C05-CH11}
% CORE-DERIVATION-PLAN: KN-V1-C05-DERIVATION-001
\SLKnowledgeBridge{同时处理内部驻点与全为 $0$、全为 $1$ 的边界样本。}{独立样本 $x_i\in\{0,1\}$、成功次数 $s=\sum_i x_i$ 与参数 $\theta\in[0,1]$。}{参数空间是闭区间；当 $s=0$ 或 $s=n$ 时不能只依赖内部一阶条件。}{先写对数似然 $\ell(\theta)=s\log\theta+(n-s)\log(1-\theta)$。}

\begin{derivationgoalbox}
推导伯努利参数的极大似然估计，并说明边界样本的处理。
\end{derivationgoalbox}
\begin{derivationprepbox}
数据 $x_i\in\{0,1\}$ 独立同分布，令 $s=\sum_ix_i$，参数空间为闭区间 $[0,1]$。
\end{derivationprepbox}
\textbf{推导第1步：}写出 $L(p)=p^s(1-p)^{n-s}$，在 $0<p<1$ 上有 $\ell(p)=s\log p+(n-s)\log(1-p)$。

\textbf{推导第2步：}求导得到 $\ell'(p)=s/p-(n-s)/(1-p)$，令其为零并通分，得到 $s(1-p)=(n-s)p$。

\textbf{推导第3步：}解得 $\hat p=s/n=\bar x$。当 $0<s<n$ 时二阶导数为负；当 $s=0$ 或 $s=n$ 时，比较闭区间边界分别得到0或1，仍与 $s/n$ 一致。

\phantomsection\label{struct:V1-C05-CH14}
% KNOWLEDGE-PLAN: KN-V1-C05-COROLLARY-001 L1
\paragraph{读前自检：Corollary：伯努利样本中，样本均值同时是成功概率 p 的极大似然估计与无偏估计。}伯努利样本令$s=\sum_i x_i$；请对$\ell(p)=s\log p+(n-s)\log(1-p)$求导并解驻点，核对$\widehat p=s/n=\bar x$且$\mathbb E[\bar X]=p$。\par

\begin{corollary}
伯努利样本中，样本均值同时是成功概率$p$的极大似然估计与无偏估计。
\end{corollary}

\phantomsection\label{struct:V1-C05-CH20}
\begin{example}[十次伯努利试验的极大似然]\label{exm:V1-C05-binomial-mle}
观察到10次试验中7次成功。求伯努利成功概率的MLE，并比较候选参数$0.5$与$0.7$的对数似然。
\end{example}
\SLExampleSolutionHeading{exm:V1-C05-binomial-mle}
\begin{solution}
\SLReadTranslation{由十次伯努利试验中的七次成功估计成功概率，并比较$p=0.7$与$p=0.5$的对数似然。}
\SolGiven 成功数$s=7$、失败数$n-s=3$且$0<s<n=10$；参数域为$p\in[0,1]$，内部对数似然可微。
\SLMethodTrigger{写出二项对数似然并解得分方程，以严格凹性和端点行为认证极大值，再直接比较两个候选。}
\SolPlan 在$(0,1)$上求解$\ell'(p)=0$，用$\ell''(p)<0$及边界似然排除其他极大点，并计算$\ell(0.7)-\ell(0.5)$。
\SolDerive
参数域为$p\in[0,1]$，成功数$s=7$满足$0<s<n=10$，故极大点位于内部。省略与$p$无关的组合常数，
\[
\begin{aligned}
\ell(p)&=7\log p+3\log(1-p),\\
\ell'(p)&=\frac7p-\frac3{1-p}=0
\Longleftrightarrow 7(1-p)=3p
\Longleftrightarrow \widehat p=0.7,\\
\ell''(p)&=-\frac7{p^2}-\frac3{(1-p)^2}<0,\\
\ell(0.7)-\ell(0.5)
&=7\log\frac{0.7}{0.5}+3\log\frac{0.3}{0.5}>0.
\end{aligned}
\]
\SolCheck 严格凹性使内部驻点唯一；成功与失败次数均为正，故$p=0,1$处似然为$0$而$p=0.7$处为正。另有$\ell(0.7)-\ell(0.5)\approx0.82283>0$。

\SolAnswer 唯一的极大似然估计为$\widehat p=0.7=7/10$；候选$0.7$的对数似然比$0.5$大约$0.82283$。
\end{solution}


% KNOWLEDGE-PLAN: KN-V1-C05-FORMULA-004 L1
\paragraph{读前自检：先验、似然与后验。}在“贝叶斯公式写成 $p(\theta\mid x)\propto p(x\mid\theta)p(\theta)$”成立时处理先验、似然与后验：请把观测对应因子代入$p(\theta\mid x)\propto p(x\mid\theta)p(\theta)$并合并一次，并以先验表达观测数据之前的不确定性，后验同时吸收先验与样本信息判定结果。\par
% KNOWLEDGE-PLAN: KN-V1-C05-FORMULA-005 L1
\paragraph{读前自检：最大后验估计。}最大后验估计在先验$p(\theta)>0$的参数域上最大化$\log p(x\mid\theta)+\log p(\theta)$；请对一个候选代入两项，核对均匀先验时排序退化为极大似然。\par
% KNOWLEDGE-PLAN: KN-V1-C05-FORMULA-006 L1
\paragraph{读前自检：共轭先验。}围绕常用分布与统计推断中的共轭先验，先采用条件“若先验与后验属于同一分布族”，再把“若先验与后验属于同一分布族”中的已声明对象代回常用分布与统计推断中的共轭先验的定义；最后验证共轭性使更新只需改变有限维超参数。\par

\SLTeachSection{先验、后验、共轭与最大后验}{sec:V1-C05-S06}
\knowledgeanchor{PRE-PR-024}{先验、似然与后验}
贝叶斯公式写成 $p(\theta\mid x)\propto p(x\mid\theta)p(\theta)$。先验表达观测数据之前的不确定性，后验同时吸收先验与样本信息。
\knowledgeanchor{PRE-PR-025}{最大后验估计}
$\hat\theta_{\mathrm{MAP}}\in\arg\max_\theta\{\log p(x\mid\theta)+\log p(\theta)\}$，可视为带先验正则项的优化。
\knowledgeanchor{PRE-PR-030}{共轭先验}
若先验与后验属于同一分布族，则称先验对该似然共轭。共轭性使更新只需改变有限维超参数。
\knowledgeanchor{PRE-PR-032}{狄利克雷分布}
令
\[
\Delta_{K-1}^{\circ}
=\{\boldsymbol p\in(0,1)^K:\textstyle\sum_{k=1}^Kp_k=1\},
\qquad \alpha_k>0,
\qquad \alpha_0=\sum_{k=1}^K\alpha_k.
\]
相对于坐标$(p_1,\ldots,p_{K-1})$上的$(K-1)$维Lebesgue参考测度，狄利克雷分布在单纯形相对内部的密度为
\[
f(\boldsymbol p\mid\boldsymbol\alpha)
=\frac{\Gamma(\alpha_0)}{\prod_{k=1}^K\Gamma(\alpha_k)}
\prod_{k=1}^Kp_k^{\alpha_k-1},
\qquad
p_K=1-\sum_{k=1}^{K-1}p_k.
\]
它不是相对于$\mathbb R^K$中$K$维Lebesgue测度的密度。靠近面$p_k=0$时，$\alpha_k<1$、$=1$、$>1$分别对应密度沿该坐标趋于无穷、保持有限非零因子、趋于零；边界本身的参考测度为零，密度极限不等于边界上的点概率。该分布是多项分布参数的共轭先验，完整构造见\cref{chap:V5-C05}。

\phantomsection\label{struct:V1-C05-CH10}
\phantomsection\label{struct:V1-C05-CH21}
\phantomsection\label{struct:V1-C05-CH22}
\phantomsection\label{struct:V1-C05-CH23}
\phantomsection\label{struct:V1-C05-CH24}
\phantomsection\label{struct:V1-C05-CH25}
\phantomsection\label{struct:V1-C05-CH26}
\phantomsection\label{struct:V1-C05-CH27}
\begin{geometrybox}
\textbf{几何解释。}指数族的自然参数位于凸参数域，对数配分函数的Hessian半正定；因此自然参数到期望参数的映射由凸函数梯度给出。高斯等密度面则是马氏距离相同的椭球。
\end{geometrybox}
\phantomsection\label{struct:V1-C05-CH17}
\begin{probabilitybox}
边缘似然把每个参数值对数据的解释能力按先验加权，自动体现“拟合好但只在狭小参数区域有效”的模型会受到体积惩罚。隐变量模型则先描述完整数据，再对未观察部分求和。
\end{probabilitybox}
\phantomsection\label{struct:V1-C05-CH18}
\begin{optimizationbox}
\textbf{优化解释。}MLE优化数据项，MAP在其上加对数先验；隐变量似然的对数套求和造成变量耦合，后续EM算法通过下界与交替优化处理。共轭模型把函数优化简化为超参数加法更新。
\end{optimizationbox}


\phantomsection\label{struct:V1-C05-CH28}
\phantomsection\label{struct:V1-C05-CH29}
% KNOWLEDGE-PLAN: KN-V1-C05-BOUNDARY-001 L1
\paragraph{读前自检：复杂度分析：先验、后验、共轭与最大后验。}先验、后验、共轭与最大后验中的“类别标签若以整数索引给出”决定可执行范围；请由$K$的总成本剥离一次迭代或一次样本因子，结果须满足预测复杂度：输出一个$K$类后验预测分布需$O(K)$。\par

\begin{complexitybox}
\textbf{复杂度假设。}类别标签若以整数索引给出，则一次计数更新为单位成本；若以$K$维稠密向量给出，则必须读取全部$K$个分量。高斯界假设使用稠密完整协方差并采用Cholesky分解；稀疏、对角或低秩协方差应按相应数据结构重新计数。
设$n$为样本数、$K$为类别数、$d$为连续变量维数。\textbf{训练复杂度：}\DirichletMultinomial{}更新在类别标签稀疏表示下累计计数为$O(n+K)$，若每条样本以$K$维稠密向量给出则为$O(nK)$；完整协方差高斯的Cholesky分解为$O(d^3)$。\textbf{预测复杂度：}输出一个$K$类后验预测分布需$O(K)$；已分解的高斯对一个样本计算密度通常需$O(d^2)$。\textbf{存储复杂度：}类别计数与参数为$O(K)$，完整协方差及其分解为$O(d^2)$。应复用分解而不是反复求逆。
\end{complexitybox}

\phantomsection\label{struct:V1-C05-CH30}
\phantomsection\label{struct:V1-C05-CH31}
\begin{applicabilitybox}
\textbf{适用条件：}MLE适合抽样模型可信且样本量足以识别参数的情形；MAP适合存在可辩护先验或正则化需求的情形；共轭更新适合需要快速顺序推断的标准模型。\textbf{局限性：}高斯对重尾和离群点敏感，MAP忽略后验宽度，边缘似然依赖先验尺度，隐变量模型还可能存在不可辨识与多个局部最优。
\end{applicabilitybox}

\phantomsection\label{struct:V1-C05-CH32}
\begin{warningbox}
典型误用包括：把似然当作关于参数的归一化概率；忘记参数决定支持集；混用Gamma的速率与尺度参数化；把协方差矩阵的行列式写成负数或零却仍使用密度公式；只解一阶方程而不检查边界与最大性；把MAP点估计当作完整贝叶斯预测。
\end{warningbox}

\phantomsection\label{struct:V1-C05-CH33}
\begin{comparisonbox}
\textbf{概念对比：}概率固定参数后衡量可能数据，似然固定数据后比较参数；MLE只使用似然，MAP再加入先验众数偏好，后验均值则平均全部参数不确定性。边缘似然积分掉参数，剖面似然最大化掉干扰参数，两种消元方式不同。参数是总体机制的未知量，统计量是样本的已知函数。
\end{comparisonbox}

% KNOWLEDGE-PLAN: KN-V1-C05-ALGORITHM_IDEA-001 L1
\paragraph{读前自检：狄利克雷--多项分布的顺序共轭更新：执行契约。}狄利克雷--多项分布的顺序共轭更新输入“类别数$K$、严格正Dirichlet先验$\boldsymbol\alpha$与有限类别观测流$y_{1:n}$”，条件“$K\ge1$，$n\in\mathbb N_0$，$\alpha_k>0$且有限，每个$y_i\in\{1,\ldots,K\}$”；请执行“每次有限正参数加一后原子提交$\widehat\alpha$”，以“恰好处理$n$个观测并归一化$K$项，或首次失败时停止”判停。\par
% KNOWLEDGE-PLAN: KN-V1-C05-ALGORITHM_IDEA-002 L2
\begin{keypointbox}[title={狄利克雷--多项分布的顺序共轭更新：五步阅读}]
\textbf{输入。}写出数据对象、固定参数与必须成立的条件。\par
\textbf{初始化。}标明主状态、计数器和第一个合法状态。\par
\textbf{核心更新。}只手算一次从旧状态到候选状态的更新，并指出何时提交。\par
\textbf{数学停止。}写出能直接核验的停止证书；预算耗尽不等同于收敛。\par
\textbf{输出。}列出返回对象、实际计数、中心状态和用于复核的诊断。
\end{keypointbox}

\begin{AlgorithmContract}{
  input={类别数$K$、严格正Dirichlet先验$\boldsymbol\alpha$与有限类别观测流$y_{1:n}$},
  output={最后合法后验参数$\widehat{\boldsymbol\alpha}$、后验均值$\widehat{\boldsymbol p}$、实际观测/归一化计数、状态与最终诊断},
  preconditions={$K\ge1$，$n\in\mathbb N_0$，$\alpha_k>0$且有限，每个$y_i\in\{1,\ldots,K\}$},
  initialization={复制$\widehat\alpha\leftarrow\alpha$，置$i_{\rm done}=p_{\rm done}=0$并保持均值候选未提交},
  loop={按输入次序逐观测形成单个分量加一候选，全部处理后扫描$K$个分量形成归一化均值},
  update={每次有限正参数加一后原子提交$\widehat\alpha$；完整均值候选核验后整体提交},
  domain={后验参数始终严格正且有限，归一化总量严格正，后验均值位于概率单纯形},
  failure={维数、先验、预算或类别编码非法返回$\mathtt{invalid\_input}$；参数加法、求和或归一化非有限返回$\mathtt{numerical\_failure}$并保留最后合法后验参数},
  stop={恰好处理$n$个观测并归一化$K$项，或首次失败时停止},
  budget={至多$n$次计数提交与$K$次归一化，确定工作上限$n+K$},
  status={$\mathtt{completed}$、$\mathtt{invalid\_input}$、$\mathtt{numerical\_failure}$},
  iterations={$i_{\rm done}$计实际提交观测数，$p_{\rm done}$计实际形成均值分量数},
  diagnostics={最终总浓度、概率和误差、零观测标志、失败观测/类别下标和阶段},
  complexity={索引类别输入为$O(n+K)$时间、$O(K)$空间；稠密独热输入需$O(nK)$读取时间}
}
\begin{algorithm}[H]
\caption{狄利克雷--多项分布的顺序共轭更新}\label{alg:V1-C05-dirichlet}
\KwIn{先验参数向量$\boldsymbol\alpha=(\alpha_1,\ldots,\alpha_K)^{\mathsf T}\in\R_{>0}^K$；类别观测流 $y_1,\ldots,y_n$}
\KwOut{最后合法$\widehat{\boldsymbol\alpha},\widehat{\boldsymbol p}$；$i_{\rm done},p_{\rm done}$；$\mathtt{status}$与诊断}
若$K,n$、维数、先验有限正性或任一类别编码非法，返回空结果、计数$0$、$\mathtt{invalid\_input}$及位置诊断\;
初始化$\widehat{\boldsymbol\alpha}\leftarrow\boldsymbol\alpha,i_{\rm done}\leftarrow0,p_{\rm done}\leftarrow0$\;
\For{$i\leftarrow1$ \KwTo $n$}{
  复制$\boldsymbol\alpha^+\leftarrow\widehat{\boldsymbol\alpha}$并令$\alpha^+_{y_i}\leftarrow\alpha^+_{y_i}+1$；若非有限则返回$\widehat{\boldsymbol\alpha}$、$\mathtt{numerical\_failure}$及$i$诊断\;
  原子提交$\widehat{\boldsymbol\alpha}\leftarrow\boldsymbol\alpha^+$并增加$i_{\rm done}$\;
}
计算$\widehat\alpha_0\leftarrow\sum_k\widehat\alpha_k$；若$\widehat\alpha_0$非有限或不正，返回$\widehat{\boldsymbol\alpha}$、$\mathtt{numerical\_failure}$及总浓度诊断\;
\For{$k\leftarrow1$ \KwTo $K$}{形成$p_k^+\leftarrow\widehat\alpha_k/\widehat\alpha_0$；若非有限则返回$\widehat{\boldsymbol\alpha}$、$\mathtt{numerical\_failure}$及$k$诊断；否则增加$p_{\rm done}$\;}
若$p^+$非负性或和为$1$核验失败，返回最后合法$\widehat\alpha$、$\mathtt{numerical\_failure}$及概率和诊断\;
原子提交$\widehat{\boldsymbol p}\leftarrow\boldsymbol p^+$；返回完整结果、$\mathtt{completed}$及总浓度/零观测最终诊断\;
\end{algorithm}
\end{AlgorithmContract}
\SLRobustImplementationStrip{首次阅读只跟踪$\widehat{\boldsymbol\alpha}=\boldsymbol\alpha+\boldsymbol n$与按$\widehat\alpha_0$归一化；实现时再检查类别编码、有限性和状态}

\textbf{终止与定义域说明。}顺序共轭更新对每个观测恰执行一次计数加法，$n$有限时必在有限步终止；它不是数值迭代，因此无需渐近收敛判据。参数非正或类别编码越界时，相应概率模型未定义，必须停止而不能继续归一化。


% KNOWLEDGE-PLAN: KN-V1-C05-FORMULA-007 L1
\paragraph{读前自检：边缘似然。}常用分布与统计推断中的边缘似然要求先确认“边缘似然又称模型证据：$p(x)=\int p(x\mid\theta)p(\theta)\,\mathrm d\theta$”；请随后把观测对应因子代入$p(x)$并合并一次，用它平均参数不确定性，而不是把参数替换为一个点估计验收。\par
% KNOWLEDGE-PLAN: KN-V1-C05-CONCEPT-012 L1
\paragraph{读前自检：隐变量与不完全数据。}联合模型$p(x,z\mid\theta)$中$z$未观测；请对$z$求和得到$p(x\mid\theta)=\sum_zp(x,z\mid\theta)$，再核对观测对数似然为$\log\sum_zp(x,z\mid\theta)$而不是$\sum_z\log p(x,z\mid\theta)$。\par
% KNOWLEDGE-PLAN: KN-V1-C05-CONCEPT-013 L1
\paragraph{读前自检：指数族分布。}指数族分布依赖“并假设存在$\eta$的开邻域$N$及可积函数$g_1,g_2$”；请据此把$\frac{\partial^2A}{\partial\eta_i\partial\eta_j}$在其支持上求和或积分一次，并直接核对对期望的第$i$个分量关于$\eta_j$求导。\par

\SLAdvancedSection{边缘似然、隐变量与指数族}{sec:V1-C05-S07}
\knowledgeanchor{PRE-PR-028}{边缘似然}
边缘似然又称模型证据：$p(x)=\int p(x\mid\theta)p(\theta)\,\mathrm d\theta$。它平均参数不确定性，而不是把参数替换为一个点估计。
\knowledgeanchor{PRE-PR-029}{隐变量与不完全数据}
若联合模型为 $p(x,z\mid\theta)$ 而 $z$ 未观测，则观测似然是 $p(x\mid\theta)=\sum_zp(x,z\mid\theta)$ 或相应积分。隐变量的边缘化常使对数似然出现“对数套求和”。
\knowledgeanchor{PRE-PR-034}{指数族分布}
指数族统一写作
\[
p(x\mid\eta)=h(x)\exp\{\eta^{\mathsf T}T(x)-A(\eta)\},
\]
其中 $T(x)$ 是充分统计量，$A(\eta)$ 是对数配分函数。

% DERIVATION-ID: V1-C05-log-partition-derivatives
% CORE-DERIVATION-PLAN: KN-V1-C05-DERIVATION-002
\SLKnowledgeBridge{由对数配分函数导数得到充分统计量的均值、协方差与凸性。}{$A(\boldsymbol\eta)=\log Z(\boldsymbol\eta)$、充分统计量 $T(X)$ 与自然参数 $\boldsymbol\eta$。}{$\boldsymbol\eta\in\operatorname{int}\mathcal H$；支配条件允许两次交换求导与积分。}{先对 $Z(\boldsymbol\eta)=\int h(x)e^{\boldsymbol\eta^{\mathsf T}T(x)}\,dx$ 逐分量求导。}

\begin{derivationgoalbox}
推导$\nabla A(\eta)$与$\nabla^2A(\eta)$，解释自然参数、期望参数和凸性之间的关系。
\end{derivationgoalbox}
\begin{derivationprepbox}
设自然参数空间$\mathcal H=\{\eta:Z(\eta)<\infty\}$，取$\eta\in\operatorname{int}\mathcal H$，$A(\eta)=\log Z(\eta)$。并假设存在$\eta$的开邻域$N$及可积函数$g_1,g_2$，使所有$\xi\in N$与几乎处处的$x$满足$|T_j(x)|h(x)e^{\xi^{\mathsf T}T(x)}\le g_1(x)$、$|T_i(x)T_j(x)|h(x)e^{\xi^{\mathsf T}T(x)}\le g_2(x)$。这些支配条件允许两次交换求导与积分；边界参数须另行处理。
\end{derivationprepbox}
\textbf{推导第1步：对归一化常数求导。}第$j$个分量满足
\[
\frac{\partial Z}{\partial\eta_j}
=\int T_j(x)h(x)e^{\eta^{\mathsf T}T(x)}\,\mathrm dx.
\]

\textbf{推导第2步：识别期望。}由$A=\log Z$和$p_\eta(x)=h(x)e^{\eta^{\mathsf T}T(x)}/Z(\eta)$，
\[
\frac{\partial A}{\partial\eta_j}
=\frac{1}{Z}\frac{\partial Z}{\partial\eta_j}
=\mathbb E_\eta[T_j(X)],
\qquad
\nabla A(\eta)=\mathbb E_\eta[T(X)].
\]

\textbf{推导第3步：再求导并中心化。}对期望的第$i$个分量关于$\eta_j$求导，商法则给出
\[
\frac{\partial^2A}{\partial\eta_i\partial\eta_j}
=\mathbb E_\eta[T_iT_j]-\mathbb E_\eta[T_i]\mathbb E_\eta[T_j].
\]
因此$\nabla^2A(\eta)=\operatorname{Cov}_\eta[T(X)]\succeq0$；协方差矩阵半正定说明$A$为凸函数。

\phantomsection\label{struct:V1-C05-CH16}
\begin{chapterexercisebox}
\phantomsection\label{struct:V1-C05-CH34}
\phantomsection\label{struct:V1-C05-CH35}
本章共21道练习。每道题干后立即给出同号解析；长解析可自然跨页，续页标题保留题号。
\end{chapterexercisebox}

\begin{SLChapterExercisePairSection}

\Needspace{6\baselineskip}
\begin{exercise}\label{ex:V1-C05-S01-01}
某分类器对单个独立样本预测正确的概率为0.8。测试10个样本时，恰好正确8个及至少正确1个的概率各是多少？
\end{exercise}
\phantomsection\label{sol:V1-C05-S01-01}
\begin{chapterexercisesolution}{ex:V1-C05-S01-01}
令$S$为10个独立样本中的正确数，则$S\sim\operatorname{Binomial}(10,0.8)$，支持集为$\{0,\ldots,10\}$，所有概率均有限。于是
\[
\begin{aligned}
P(S=8)
&=\binom{10}{8}(0.8)^8(0.2)^2
=0.301989888,\\
P(S\ge1)
&=1-P(S=0)
=1-(0.2)^{10}
=0.9999998976.
\end{aligned}
\]
第二式使用支持集中的唯一补事件$S=0$，避免逐项累加。
\end{chapterexercisesolution}

\Needspace{6\baselineskip}
\begin{exercise}\label{ex:V1-C05-S01-02}
证明多项分布中 $E[N_k]=np_k$，并求 $\operatorname{Cov}(N_j,N_k)$（$j\ne k$）。
\end{exercise}
\phantomsection\label{sol:V1-C05-S01-02}
\begin{chapterexercisesolution}{ex:V1-C05-S01-02}
设$n\in\mathbb N$，类别概率$p_k\ge0$且$\sum_kp_k=1$。令$I_{ik}=\mathbf1\{Z_i=k\}$，则$I_{ik}$有界，所需期望均存在，且
\[
\begin{aligned}
N_k&=\sum_{i=1}^n I_{ik},
&\mathbb E[N_k]&=\sum_{i=1}^n\mathbb E[I_{ik}]=np_k,\\
\operatorname{Cov}(I_{ij},I_{ik})
&=\mathbb E[I_{ij}I_{ik}]-p_jp_k
=-p_jp_k,\qquad j\ne k,\\
\operatorname{Cov}(N_j,N_k)
&=\sum_{i=1}^n\operatorname{Cov}(I_{ij},I_{ik})
=-np_jp_k.
\end{aligned}
\]
最后一步使用不同试验间独立，使跨试验协方差为零。
\end{chapterexercisesolution}

\Needspace{6\baselineskip}
\begin{exercise}\label{ex:V1-C05-S01-03}
若 $S\sim\operatorname{Binomial}(n,p)$，用指示变量证明 $E[S(S-1)]=n(n-1)p^2$。
\end{exercise}
\phantomsection\label{sol:V1-C05-S01-03}
\begin{chapterexercisesolution}{ex:V1-C05-S01-03}
令$X_i\overset{\mathrm{iid}}{\sim}\operatorname{Bernoulli}(p)$、$0\le p\le1$，且$S=\sum_{i=1}^nX_i$。因变量有界，二阶矩存在。利用$X_i^2=X_i$，
\[
\begin{aligned}
S(S-1)
&=\left(\sum_iX_i\right)^2-\sum_iX_i\\
&=\sum_i(X_i^2-X_i)+\sum_{i\ne j}X_iX_j
=\sum_{i\ne j}X_iX_j,\\
\mathbb E[S(S-1)]
&=\sum_{i\ne j}\mathbb E[X_i]\mathbb E[X_j]
=n(n-1)p^2.
\end{aligned}
\]
因此二项变量的二阶阶乘矩为$n(n-1)p^2$，边界$p=0,1$也包含在内。
\end{chapterexercisesolution}

\Needspace{6\baselineskip}
\begin{exercise}\label{ex:V1-C05-S02-01}
若 $X\sim\mathcal N(2,9)$，求 $Y=3X-1$ 的分布，并把 $P(X\le5)$ 写成标准高斯分布函数 $\Phi$。
\end{exercise}
\phantomsection\label{sol:V1-C05-S02-01}
\begin{chapterexercisesolution}{ex:V1-C05-S02-01}
$X\sim\mathcal N(2,9)$的支持集为$\mathbb R$且二阶矩有限。对$Y=3X-1$，
\[
\begin{aligned}
\mathbb E[Y]&=3\mathbb E[X]-1=3\cdot2-1=5,\\
\operatorname{Var}(Y)&=3^2\operatorname{Var}(X)=9\cdot9=81,\\
Y&\sim\mathcal N(5,81),\\
P(X\le5)
&=P\left(\frac{X-2}{3}\le\frac{5-2}{3}\right)=\Phi(1).
\end{aligned}
\]
所以$Y\sim\mathcal N(5,81)$，且$P(X\le5)=\Phi(1)$。
\end{chapterexercisesolution}

\Needspace{6\baselineskip}
\begin{exercise}\label{ex:V1-C05-S02-02}
设 $X\sim\mathcal N(\mu,\Sigma)$，求标量 $a^{\mathsf T}X+b$ 的均值与方差。
\end{exercise}
\phantomsection\label{sol:V1-C05-S02-02}
\begin{chapterexercisesolution}{ex:V1-C05-S02-02}
声明$X\in\mathbb R^d$、$\mu,a\in\mathbb R^d$、$\Sigma\in\mathbb R^{d\times d}$对称半正定、$b\in\mathbb R$。高斯二阶矩存在，故
\[
\begin{aligned}
\mathbb E[a^{\mathsf T}X+b]&=a^{\mathsf T}\mu+b,\\
\operatorname{Var}(a^{\mathsf T}X+b)
&=\mathbb E\!\left[a^{\mathsf T}(X-\mu)(X-\mu)^{\mathsf T}a\right]\\
&=a^{\mathsf T}\Sigma a\ge0,\\
a^{\mathsf T}X+b
&\sim\mathcal N(a^{\mathsf T}\mu+b,\ a^{\mathsf T}\Sigma a).
\end{aligned}
\]
若$a^{\mathsf T}\Sigma a=0$，最后一项表示退化在其均值处的点质量。
\end{chapterexercisesolution}

\Needspace{6\baselineskip}
\begin{exercise}\label{ex:V1-C05-S02-03}
二维高斯协方差为 $\Sigma=\begin{psmallmatrix}4&2\\2&9\end{psmallmatrix}$。求两分量相关系数，并检验矩阵正定。
\end{exercise}
\phantomsection\label{sol:V1-C05-S02-03}
\begin{chapterexercisesolution}{ex:V1-C05-S02-03}
协方差矩阵$\Sigma\in\mathbb R^{2\times2}$对称，边缘方差$4,9$均为正，因此相关系数有定义。计算
\[
\begin{aligned}
\sigma_1&=\sqrt4=2,\qquad \sigma_2=\sqrt9=3,\\
\rho_{12}&=\frac{\Sigma_{12}}{\sigma_1\sigma_2}
=\frac2{2\cdot3}=\frac13,\\
\Delta_1&=4>0,\qquad
\Delta_2=\det(\Sigma)=4\cdot9-2^2=32>0.
\end{aligned}
\]
由Sylvester判据$\Sigma\succ0$，故该协方差合法且非退化。
\end{chapterexercisesolution}

\Needspace{6\baselineskip}
\begin{exercise}\label{ex:V1-C05-S03-01}
用Gamma递推计算 $\Gamma(5)$，并说明它与 $5!$ 的区别。
\end{exercise}
\phantomsection\label{sol:V1-C05-S03-01}
\begin{chapterexercisesolution}{ex:V1-C05-S03-01}
Gamma函数在$a>0$时由正半轴上的有限积分定义，并满足$\Gamma(a+1)=a\Gamma(a)$、$\Gamma(1)=1$。因此
\[
\begin{aligned}
\Gamma(5)
&=4\Gamma(4)=4\cdot3\Gamma(3)
=4\cdot3\cdot2\Gamma(2)\\
&=4\cdot3\cdot2\cdot1\Gamma(1)=4!=24,\\
5!&=5\cdot4!=120.
\end{aligned}
\]
整数$n\ge1$时$\Gamma(n)=(n-1)!$，故同一自变量下不能直接写成$n!$。
\end{chapterexercisesolution}

\Needspace{6\baselineskip}
\begin{exercise}\label{ex:V1-C05-S03-02}
若 $X\sim\operatorname{Gamma}(a,b)$ 采用形状--速率参数化，求 $Y=cX$（$c>0$）的Gamma参数。
\end{exercise}
\phantomsection\label{sol:V1-C05-S03-02}
\begin{chapterexercisesolution}{ex:V1-C05-S03-02}
采用形状—速率参数化，要求$a,b,c>0$；$X$与$Y=cX$的支持集均为$(0,\infty)$。对$y>0$作变量代换$x=y/c$，
\[
\begin{aligned}
f_Y(y)
&=f_X(y/c)\left|\frac{\mathrm d(y/c)}{\mathrm dy}\right|\\
&=\frac{b^a}{\Gamma(a)}\left(\frac yc\right)^{a-1}
e^{-by/c}\frac1c\\
&=\frac{(b/c)^a}{\Gamma(a)}y^{a-1}e^{-(b/c)y}.
\end{aligned}
\]
故$Y\sim\operatorname{Gamma}(a,b/c)$，并有$\mathbb E[Y]=c(a/b)=a/(b/c)$。
\end{chapterexercisesolution}

\Needspace{6\baselineskip}
\begin{exercise}\label{ex:V1-C05-S03-03}
设 $p=(p_1,p_2,p_3)$ 服从参数 $(2,3,5)$ 的狄利克雷分布，求各分量期望，并解释它们为何相加为1。
\end{exercise}
\phantomsection\label{sol:V1-C05-S03-03}
\begin{chapterexercisesolution}{ex:V1-C05-S03-03}
参数$(2,3,5)$均为正，支持集是概率单纯形
$\Delta_2=\{p_k\ge0:\sum_{k=1}^3p_k=1\}$，分量有界因而期望有限。令$\alpha_0=2+3+5=10$，则
\[
\begin{aligned}
\mathbb E[p_1]&=\frac2{10}=0.2,\qquad
\mathbb E[p_2]=\frac3{10}=0.3,\\
\mathbb E[p_3]&=\frac5{10}=0.5,\\
\sum_{k=1}^3\mathbb E[p_k]
&=\mathbb E\!\left[\sum_{k=1}^3p_k\right]
=\mathbb E[1]=1.
\end{aligned}
\]
因此三个分量的期望为$(0.2,0.3,0.5)$，并保持单纯形上的和为$1$。
\end{chapterexercisesolution}

\Needspace{6\baselineskip}
\begin{exercise}\label{ex:V1-C05-S04-01}
判断 $\bar X$、$\sum_i(X_i-\mu)^2$、$\max_iX_i$ 中哪些在未知 $\mu$ 时是统计量。
\end{exercise}
\phantomsection\label{sol:V1-C05-S04-01}
\begin{chapterexercisesolution}{ex:V1-C05-S04-01}
统计量必须是仅依赖样本$X_{1:n}$、不含未知参数的可测函数。按该判据，
\[
\begin{aligned}
T_1(X_{1:n})&=\bar X=\frac1n\sum_{i=1}^nX_i
&&\text{是统计量},\\
T_2(X_{1:n};\mu)&=\sum_{i=1}^n(X_i-\mu)^2
&&\text{在未知$\mu$时不是统计量},\\
T_3(X_{1:n})&=\max_{1\le i\le n}X_i
&&\text{是统计量}.
\end{aligned}
\]
若将$\mu$换为样本函数$\bar X$，则$\sum_i(X_i-\bar X)^2$重新成为统计量。
\end{chapterexercisesolution}

\Needspace{6\baselineskip}
\begin{exercise}\label{ex:V1-C05-S04-02}
设 $X_1,\ldots,X_n$ 独立同分布，均值为$\mu$、方差为$\sigma^2$。证明 $\bar X$ 无偏并给出标准误。
\end{exercise}
\phantomsection\label{sol:V1-C05-S04-02}
\begin{chapterexercisesolution}{ex:V1-C05-S04-02}
设$n\ge1$，$X_i$独立同分布且$\mathbb E[X_i]=\mu$、$\operatorname{Var}(X_i)=\sigma^2<\infty$，故样本均值的一、二阶矩存在。于是
\[
\begin{aligned}
\mathbb E[\bar X]
&=\frac1n\sum_{i=1}^n\mathbb E[X_i]
=\frac{n\mu}{n}=\mu,\\
\operatorname{Var}(\bar X)
&=\frac1{n^2}\sum_{i=1}^n\operatorname{Var}(X_i)
=\frac{\sigma^2}{n},\\
\operatorname{SE}(\bar X)
&=\sqrt{\operatorname{Var}(\bar X)}
=\frac{\sigma}{\sqrt n}.
\end{aligned}
\]
第一式证明无偏，独立性只在第二式消去交叉协方差时使用。
\end{chapterexercisesolution}

\Needspace{6\baselineskip}
\begin{exercise}\label{ex:V1-C05-S04-03}
解释为什么观测到数据后样本值是固定的，但频率学派仍把估计量称为随机变量。
\end{exercise}
\phantomsection\label{sol:V1-C05-S04-03}
\begin{chapterexercisesolution}{ex:V1-C05-S04-03}
对象分为抽样前的随机样本$X_{1:n}$与观测后的实现$x_{1:n}$。同一估计规则$T$在两层分别写成
\[
\begin{aligned}
\widehat\theta&=T(X_{1:n})
&&\text{是随机变量},\\
\widehat\theta_{\mathrm{obs}}&=T(x_{1:n})
&&\text{是固定数值},\\
\operatorname{Bias}_\theta(T)
&=\mathbb E_\theta[T(X_{1:n})]-\theta,\qquad
\operatorname{Var}_\theta(T)=\mathbb E_\theta[(T-\mathbb E_\theta T)^2].
\end{aligned}
\]
只要相应期望存在，偏差和方差描述重复抽样分布，而不表示已经记录的数据仍在变化。
\end{chapterexercisesolution}

\Needspace{6\baselineskip}
\begin{exercise}\label{ex:V1-C05-S05-01}
独立高斯样本方差 $\sigma^2$ 已知，推导均值 $\mu$ 的极大似然估计。
\end{exercise}
\phantomsection\label{sol:V1-C05-S05-01}
\begin{chapterexercisesolution}{ex:V1-C05-S05-01}
设$n\ge1$、$x_i\in\mathbb R$、已知$\sigma^2>0$，参数域为$\mu\in\mathbb R$；高斯密度在全支持上严格为正。忽略与$\mu$无关的常数，
\[
\begin{aligned}
\ell(\mu)
&=-\frac1{2\sigma^2}\sum_{i=1}^n(x_i-\mu)^2+C,\\
\ell'(\mu)
&=\frac1{\sigma^2}\sum_{i=1}^n(x_i-\mu)=0
\Longleftrightarrow \widehat\mu=\bar x,\\
\ell''(\mu)&=-\frac n{\sigma^2}<0.
\end{aligned}
\]
对数似然严格凹且参数域无边界，故$\widehat\mu=\bar x$是唯一全局MLE。
\end{chapterexercisesolution}

\Needspace{6\baselineskip}
\begin{exercise}\label{ex:V1-C05-S05-02}
说明对数似然为何比原似然更适合数值计算，并指出它何时不能定义。
\end{exercise}
\phantomsection\label{sol:V1-C05-S05-02}
\begin{chapterexercisesolution}{ex:V1-C05-S05-02}
对独立样本定义$L(\theta)=\prod_{i=1}^np_\theta(x_i)$及正似然域$\Theta_+=\{\theta:L(\theta)>0\}$。在$\Theta_+$上，
\[
\begin{aligned}
\ell(\theta)
&=\log L(\theta)
=\sum_{i=1}^n\log p_\theta(x_i),\\
\arg\max_{\theta\in\Theta_+}L(\theta)
&=\arg\max_{\theta\in\Theta_+}\ell(\theta),\\
L(\theta)=0
&\Longrightarrow \log L(\theta)\text{在普通实数中无定义},
\quad \ell(\theta):=-\infty\text{按扩展实数约定}.
\end{aligned}
\]
求和避免大量小数连乘下溢；零概率边界表示观测在该模型下不可能，不能用数值常数随意改写。
\end{chapterexercisesolution}

\Needspace{6\baselineskip}
\begin{exercise}\label{ex:V1-C05-S05-03}
均匀分布 $X_i\sim\operatorname{Unif}(0,\theta)$，$\theta>0$。求MLE并解释为何不能只令普通导数为零。
\end{exercise}
\phantomsection\label{sol:V1-C05-S05-03}
\begin{chapterexercisesolution}{ex:V1-C05-S05-03}
参数域为$\Theta=(0,\infty)$，单样本支持为$[0,\theta]$。令$x_{(1)}=\min_i x_i$、$x_{(n)}=\max_i x_i$，完整似然为
\[
\begin{aligned}
L(\theta;x_{1:n})
&=\theta^{-n}\mathbf1\{x_{(1)}\ge0\}
\mathbf1\{\theta\ge x_{(n)}\},\qquad \theta>0.
\end{aligned}
\]
若样本与模型相容且$x_{(n)}>0$，可行参数集为$[x_{(n)},\infty)$，并且
\[
\begin{aligned}
\frac{\mathrm d}{\mathrm d\theta}\theta^{-n}
&=-n\theta^{-n-1}<0,\qquad \theta\ge x_{(n)},\\
L(\theta)\text{严格递减}
&\Longrightarrow \widehat\theta=x_{(n)}.
\end{aligned}
\]
若$x_{(1)}<0$，则$L\equiv0$，原始似然的$\arg\max$退化为整个$\Theta$而对数似然处处为$-\infty$；若所有$x_i=0$，则$L(\theta)=\theta^{-n}\to\infty$当$\theta\downarrow0$，只有上确界而无MLE。参数改变支持集，因此内部驻点法不适用。
\end{chapterexercisesolution}

\Needspace{6\baselineskip}
\begin{exercise}\label{ex:V1-C05-S06-01}
伯努利参数先验为 $\operatorname{Beta}(2,3)$，观察8次成功、2次失败。求后验参数与后验均值。
\end{exercise}
\phantomsection\label{sol:V1-C05-S06-01}
\begin{chapterexercisesolution}{ex:V1-C05-S06-01}
参数$p$的支持为$(0,1)$，先验$\operatorname{Beta}(2,3)$正规且10次观测含8次成功、2次失败。后验核为
\[
\begin{aligned}
\pi(p\mid x)
&\propto p^{2-1}(1-p)^{3-1}\,p^8(1-p)^2\\
&=p^{10-1}(1-p)^{5-1},\\
p\mid x&\sim\operatorname{Beta}(10,5),\\
\mathbb E[p\mid x]&=\frac{10}{10+5}=\frac23.
\end{aligned}
\]
后验形状参数均为正，故归一化常数与均值都有限。
\end{chapterexercisesolution}

\Needspace{6\baselineskip}
\begin{exercise}\label{ex:V1-C05-S06-02}
对上题求MAP，并说明形状参数不大于1时为何要检查边界。
\end{exercise}
\phantomsection\label{sol:V1-C05-S06-02}
\begin{chapterexercisesolution}{ex:V1-C05-S06-02}
后验为$\operatorname{Beta}(a,b)$且本题$(a,b)=(10,5)$，两者均大于1，故密度在$(0,1)$有限且边界趋零。内部对数密度满足
\[
\begin{aligned}
\ell(p)&=(a-1)\log p+(b-1)\log(1-p)+C,\\
\ell'(p)&=\frac{a-1}{p}-\frac{b-1}{1-p}=0,\\
p_{\mathrm{MAP}}
&=\frac{a-1}{a+b-2}
=\frac9{13},\\
\ell''(p)&=-\frac{a-1}{p^2}-\frac{b-1}{(1-p)^2}<0.
\end{aligned}
\]
若$a\le1$或$b\le1$，端点可能取到最大、出现平台或使密度发散，内部公式不再给出全局MAP，必须按支持边界另行比较。
\end{chapterexercisesolution}

\Needspace{6\baselineskip}
\begin{exercise}\label{ex:V1-C05-S06-03}
三类先验为Dirichlet$(1,1,1)$，观测计数为 $(4,2,0)$。求后验与后验预测概率。
\end{exercise}
\phantomsection\label{sol:V1-C05-S06-03}
\begin{chapterexercisesolution}{ex:V1-C05-S06-03}
类别概率$\boldsymbol p$位于三类单纯形，先验参数均为正，计数为$\boldsymbol n=(4,2,0)$且样本数$n=6$。共轭更新为
\[
\begin{aligned}
\boldsymbol\alpha'&=(1,1,1)+(4,2,0)=(5,3,1),\\
\boldsymbol p\mid x&\sim\operatorname{Dirichlet}(5,3,1),\\
P(X_{\mathrm{new}}=k\mid x)
&=\mathbb E[p_k\mid x]
=\frac{\alpha_k'}{\sum_j\alpha_j'},\\
\boldsymbol P_{\mathrm{pred}}
&=\left(\frac59,\frac39,\frac19\right).
\end{aligned}
\]
三个预测概率非负且和为1；第三类因正先验伪计数而未被赋零概率。
\end{chapterexercisesolution}

\Needspace{6\baselineskip}
\begin{exercise}\label{ex:V1-C05-S07-01}
模型有两个参数值，先验各为 $1/2$，数据似然分别为0.1和0.4。求边缘似然与两个参数的后验概率。
\end{exercise}
\phantomsection\label{sol:V1-C05-S07-01}
\begin{chapterexercisesolution}{ex:V1-C05-S07-01}
参数支持集为$\{\theta_1,\theta_2\}$，先验质量各为$1/2$；两个似然有限且证据概率将严格为正。由全概率与Bayes公式，
\[
\begin{aligned}
p(x)&=\sum_{j=1}^2p(x\mid\theta_j)\pi_j
=0.5(0.1)+0.5(0.4)=0.25,\\
P(\theta_1\mid x)&=\frac{0.05}{0.25}=0.2,\\
P(\theta_2\mid x)&=\frac{0.20}{0.25}=0.8,\\
P(\theta_1\mid x)+P(\theta_2\mid x)&=1.
\end{aligned}
\]
所以后验分布为$(0.2,0.8)$，归一化检查通过。
\end{chapterexercisesolution}

\Needspace{6\baselineskip}
\begin{exercise}\label{ex:V1-C05-S07-02}
混合模型 $p(x)=\sum_{k=1}^K\pi_kp(x\mid z=k)$ 中，写出给定$x$后隐类别的后验责任度。
\end{exercise}
\phantomsection\label{sol:V1-C05-S07-02}
\begin{chapterexercisesolution}{ex:V1-C05-S07-02}
类别支持为$\{1,\ldots,K\}$，权重$\pi_k\ge0$且$\sum_k\pi_k=1$。对给定$x$定义有限证据
$Z(x)=\sum_{j=1}^K\pi_jp(x\mid z=j)$；仅当$0<Z(x)<\infty$时责任度有定义：
\[
\begin{aligned}
\gamma_k(x)
&=P(z=k\mid x)
=\frac{\pi_kp(x\mid z=k)}{Z(x)},\\
\gamma_k(x)&\ge0,\qquad
\sum_{k=1}^K\gamma_k(x)
=\frac{Z(x)}{Z(x)}=1,\\
\gamma_k(x)&=\mathbb E[\mathbf1\{z=k\}\mid x].
\end{aligned}
\]
若$Z(x)=0$，观测落在混合分布支持之外，后验责任度未定义。
\end{chapterexercisesolution}

\Needspace{6\baselineskip}
\begin{exercise}\label{ex:V1-C05-S07-03}
证明指数族中 $\nabla^2A(\eta)=\operatorname{Cov}_\eta[T(X)]$，并说明 $A$ 的凸性。
\end{exercise}
\phantomsection\label{sol:V1-C05-S07-03}
\begin{chapterexercisesolution}{ex:V1-C05-S07-03}
设自然参数$\eta$位于开凸集$\mathcal H$，配分积分$Z(\eta)=\int h(x)e^{\eta^{\mathsf T}T(x)}\,\mathrm dx$有限且正，并存在可积支配函数允许对$\eta$作两次积分号下求导；同时$\mathbb E_\eta\|T(X)\|_2^2<\infty$。令$A(\eta)=\log Z(\eta)$，则
\[
\begin{aligned}
\nabla A(\eta)
&=\frac{\int T(x)h(x)e^{\eta^{\mathsf T}T(x)}\,\mathrm dx}{Z(\eta)}
=\mathbb E_\eta[T(X)],\\
\nabla^2A(\eta)
&=\mathbb E_\eta[T(X)T(X)^{\mathsf T}]
-\mathbb E_\eta[T(X)]\mathbb E_\eta[T(X)]^{\mathsf T}\\
&=\operatorname{Cov}_\eta[T(X)].
\end{aligned}
\]
对任意$v$，$v^{\mathsf T}\nabla^2A(\eta)v=\operatorname{Var}_\eta(v^{\mathsf T}T(X))\ge0$，故$A$在$\mathcal H$上凸。若缺少换导与二阶矩条件，上述形式推导不能直接作为证明。
\end{chapterexercisesolution}

\end{SLChapterExercisePairSection}

\Needspace{4\baselineskip}
% KNOWLEDGE-PLAN: KN-V1-C05-CONCEPT-014 L1
\paragraph{读前自检：本章结论与下一步。}常用分布与统计推断章末由“常用分布、参数、统计量与充分统计量”落到“先核对支持边界与换导条件，再求驻点并比较边界”；请把前者产出和后者输入逐项对齐，固定核对前者末项的输出类型能否作为后者首项的输入类型。\par

\begin{SLCompactClosure}
\phantomsection\label{struct:V1-C05-CH36}
\SLChapterFiveSummary{常用分布、参数、统计量与充分统计量}{似然/MAP优化、共轭更新及配分函数求导}{先核对支持边界与换导条件，再求驻点并比较边界}{参数估计、隐变量建模和概率预测}{进入\cref{chap:V1-C06}，用KL与交叉熵比较分布}

\phantomsection\label{struct:V1-C05-CH37}
\begin{selfcheckbox}
\begin{enumerate}[leftmargin=2.2em]
  \item 能否在写密度前先声明支持集与参数化，并检查归一化？
  \item 能否从独立抽样模型逐步写出似然、对数似然、一阶条件及边界检查？
  \item 能否说明共轭先验的超参数如何随充分统计量更新？
  \item 能否区分边缘似然、后验、MAP与后验预测各自积分或优化的对象？未通过时依次回看\cref{sec:V1-C05-S01,sec:V1-C05-S02,sec:V1-C05-S03,sec:V1-C05-S05,sec:V1-C05-S06,sec:V1-C05-S07}。
\end{enumerate}
\end{selfcheckbox}
\end{SLCompactClosure}
